\chapter*{Appendix U: The Type-Theoretic Kernel}
\addcontentsline{toc}{chapter}{Appendix U: The Type-Theoretic Kernel}

\section{Overview}

This appendix presents the dependent type-theoretic core supporting the formal verification of RSVP field dynamics, Polyxan rewriting, reconciliation operators, and embedding structures. The kernel defines the admissible types, formation rules, modal operators, and inductive constructions compatible with curvature-, entropy-, and normalisation constraints.

\section{Judgemental Structure}

\subsection{Typing Judgements}

Typing judgements have the standard form:
\[
\Gamma \vdash t : A,
\]
where $\Gamma$ is a context of variables annotated with types stable under the modal invariants.

\subsection{Context Formation}

A context is a finite sequence:
\[
\Gamma = x_1 : A_1, \ldots, x_n : A_n,
\]
where each $A_i$ satisfies:
\begin{enumerate}
\item curvature coherence,
\item entropy monotonicity,
\item quine stability,
\item reconciliation compatibility.
\end{enumerate}

\section{Base Types}

\subsection{Field Types}

Field types include:
\[
\mathsf{Scalar}, \qquad 
\mathsf{Vector}, \qquad 
\mathsf{Entropy},
\]
representing components of the RSVP fields $(\Phi, \mathbf{v}, S)$.

Each carries a metric structure inherited from the embedding space.

\subsection{Graph Types}

Hypergraph types include:
\[
\mathsf{Vertex}, \qquad 
\mathsf{Edge}, \qquad
\mathsf{Cell}, \qquad
\mathsf{Tag}.
\]

These types support inductive operations compatible with the Polyxan rewriting semantics.

\subsection{Embedding Types}

Embedding types classify maps:
\[
\iota : \mathcal{G} \to M,
\]
and carry dependent structure expressing curvature-matching and harmonicity constraints.

\section{Dependent Types}

\subsection{Families of Field Configurations}

A dependent type of field configurations is given by:
\[
\mathsf{FieldConf}(n) := \{\, f : \mathbb{R}^n \to \mathbb{R} \mid f \text{ satisfies harmonicity and bounded curvature}\,\}.
\]

\subsection{Families of Hypergraph Incidences}

A dependent family of incidences is:
\[
\mathsf{Inc}(v) := \{\, c \in \mathrm{Inc}(\mathcal{G}) \mid v \in c \,\}.
\]

\subsection{Dependent Embedding Families}

For each cell $c$:
\[
\mathsf{Embed}(c) := \{\, \iota(c) \in M \mid \mathcal{R}(\iota(c)) = \mathcal{H}_0(c) \,\}.
\]

\section{Modal Types}

The kernel includes two modal operators corresponding to the verification predicates:
\[
\Box_{\mathrm{RSVP}}, \qquad 
\Box_{\mathrm{Polyx}}.
\]

\subsection{Formation Rules}

If $\Gamma \vdash A : \mathsf{Type}$ and $A$ satisfies the appropriate invariants, then:
\[
\Gamma \vdash \Box_{\mathrm{RSVP}} A : \mathsf{Type}, \qquad
\Gamma \vdash \Box_{\mathrm{Polyx}} A : \mathsf{Type}.
\]

\subsection{Introduction Rules}

A term $t : A$ yields:
\[
\Gamma \vdash t^{\ast} : \Box_{\mathrm{RSVP}} A
\]
when $t$ preserves Lyapunov energy and harmonicity.

Similarly for $\Box_{\mathrm{Polyx}}$ when $t$ preserves rewriting invariants.

\subsection{Elimination Rules}

If $\Gamma \vdash u : \Box_{\mathrm{RSVP}} A$ and $\Gamma, x : A \vdash B(x)$, then:
\[
\Gamma \vdash \mathsf{let}\; x \leftarrow u \; \mathsf{in}\; b : B(x),
\]
provided $B$ is RSVP-invariant.

Identical rules hold for $\Box_{\mathrm{Polyx}}$.

\section{Inductive Types}

\subsection{Atomic Patches}

An inductive family of atomic patches:
\[
\mathsf{AtPatch} : \mathsf{Type}
\]
is generated by constructors for tag modification, incidence extension, and local rewriting.

\subsection{Composite Patches}

Composite patches are defined inductively by:
\begin{align*}
\mathsf{CompPatch} &:=
  \mathsf{unit}
  \;\mid\;
  \mathsf{step}(p : \mathsf{AtPatch})
  \;\mid\;
  \mathsf{join}(p,q : \mathsf{CompPatch})
\end{align*}
with reconciliation forming a definitional equality on $\mathsf{join}$.

\section{Equality and Normalisation}

\subsection{Definitional Equality}

Terms $t$ and $u$ satisfy definitional equality when their embeddings yield identical curvature and entropy structure:
\[
t \equiv u \quad \text{iff} \quad
\mathcal{E}_{\mathrm{total}}[t] = \mathcal{E}_{\mathrm{total}}[u].
\]

\subsection{Normalisation}

The kernel includes a normalisation operator:
\[
\mathsf{nf} : \mathsf{CompPatch} \to \mathsf{CompPatch},
\]
satisfying idempotence and confluence.

\section{Categorical Semantics}

\subsection{Base Category}

The categorical semantics is given by a fibration:
\[
\mathcal{E} \to \mathcal{B},
\]
where $\mathcal{B}$ is the category of contexts and $\mathcal{E}$ the category of typed terms.

\subsection{Modal Semantics}

Modal types correspond to reflective subcategories:
\[
\mathbf{RSVP} \hookrightarrow \mathcal{E}, \qquad
\mathbf{Polyx} \hookrightarrow \mathcal{E}.
\]

\subsection{Inductive Type Semantics}

Inductive types correspond to initial algebras of polynomial endofunctors, with reconciliation acting as a monad multiplication.

\section{Operational Semantics}

Evaluation rules are given by a deterministic relation:
\[
t \longrightarrow t',
\]
subject to:
\begin{enumerate}
\item preservation of modal types,
\item non-increase of curvature,
\item commutation with normalisation.
\end{enumerate}

\section{Soundness and Completeness}

\subsection{Soundness}

Each typing judgement $\Gamma \vdash t : A$ implies that the semantic interpretation of $t$ lies in the correct fibre of the fibration.

\subsection{Completeness}

Every semantically valid morphism between fibres corresponds to a term of the appropriate type, up to definitional equality.

\section{Summary}

The kernel provides:
\begin{itemize}
\item dependent types for field, graph, and embedding structures,
\item modal verification types enforcing RSVP and Polyxan invariants,
\item inductive types for patches,
\item normalisation as a definitional equality scheme,
\item categorical semantics ensuring soundness.
\end{itemize}

This infrastructure supports all formal verification steps across the monograph.

